\newcommand{\docNome}{ Studio di fattibilità }   % INSERIRE LA DATA IN FORMATO aaaa-mm-gg
\newcommand{\docVersione}{0.0.1}                     % INSERIRE VERSIONE IN FORMATO x.y.z
\newcommand{\docStatus}{in lavorazione}         % AGGIORNARE SOLO QUANDO APPROVATO
\newcommand{\docUso}{}                          % INTERNO O ESTERNO
\newcommand{\docDestinatari}{} %Se verbale esterno togliere questo comando (verranno inseriti direttamente come testo) 
\newcommand{\docNomeTeam}{Sweven Team}
\newcommand{\docRedattori}{Mattia Episcopo}                    % REDATTORE
\newcommand{\docVerificatori}{}                 % VERIFICATORE
\newcommand{\docApprovazione}{}                 % APPROVATORE
\documentclass[12pt, a4paper,table]{article}
\usepackage[utf8]{inputenc}
\usepackage{fancyhdr}
\usepackage{geometry}
\usepackage{xcolor}
\usepackage{array}
\usepackage{graphicx}
\usepackage{hyperref}
\hypersetup{
pdfborder = {0 0 0}
}
\geometry{a4paper,top=3cm,bottom=3cm,left=2cm,right=2cm}
\title{\textsc{\docNome}}
\author{}
\date{}
\definecolor{footer-gray}{HTML}{808080}
\pagestyle{fancy}
\fancyhf{}
\rhead{\textcolor{footer-gray}{\docNome} }
\lhead{\textcolor{footer-gray}{Sweven Team}}
\fancyfoot{}
\cfoot{\textcolor{footer-gray}{Pagina \thepage } }
\renewcommand*\contentsname{Indice}
\begin{document}
	\maketitle
	\vspace{-3em}
	\begin{center}
	\includegraphics[scale=0.50]{images/logo.jpg} \\
	\vspace{2em}
	\huge \textsc{\docNomeTeam}\\
	\normalsize \href{mailto:swe7.team@gmail.com}{swe7.team@gmail.com}\\
	\vspace{2em}
	\begin{tabular}{r|l}
		\multicolumn{2}{c}{ \textsc{Informazioni sul documento} } \\
		\hline
		\textbf{Versione}     & \docVersione\\
		\textbf{Uso}          & \docUso\\
        	\textbf{Destinatari}  & \docDestinatari\\ %Se verbale esterno scrivere direttamente qui la lista con anche azienda Imola Informatica e professori
		\textbf{Stato}        & \docStatus\\
		\textbf{Redattori}    & \docRedattori\\
		\textbf{Verificatori} & \docVerificatori\\
		\textbf{Approvatori} & \docApprovazione\\
	\end{tabular}
	\end{center}
    \vspace{3em}
    \begin{center}
        \LARGE{\textbf{Sintesi}} 
    \end{center}
    \normalsize{Studio di fattibilità dei capitolati da scegliere}
	\thispagestyle{empty}   
	\newpage
	
	\section*{Diario delle modifiche}
	\begin{center}
	\renewcommand{\arraystretch}{1.8} %aumento ampiezza righe
	\begin{tabular}{ |c|c|m{12em}|m{7em}|c| }
	\hline
	\textbf{Versione} & \textbf{Data} & \textbf{Descrizione} &  \textbf{Autore} &  \textbf{Ruolo} \\
	\hline
	1.0.0 & aaaa-mm-gg &  & \docApprovazione & Responsabile\\ % Se il nome dell'approvatore non ci sta, metterlo a mano con aggiunta di \newline (esempio: Nome \newline Cognome)
	\hline
	0.1.0 & aaaa-mm-gg &  & \docVerificatori & Verificatore\\ % Stesso discorso di prima
	\hline
    	0.0.2 & 2022-03-25 & Aggiunto capitolato 3 & Mattia \newline Episcopo & Amministratore\\  % Stesso discorso di prima
    	0.0.1 & 2022-03-24 & Creazione iniziale del documento & Mattia \newline Episcopo & Amministratore\\  % Stesso discorso di prima
	\hline
	\end{tabular}
	\end{center}
	\newpage
	
	\tableofcontents
	\newpage
	\section{Introduzione}
		\subsection {Scopo del documento}
		Lo scopo del documento è quello di descrivere ciò che propone ogni capitolato, specificando gli aspetti positivi e le criticità secondo l'analisi del gruppo.
		\newline Sono stati esaminati solo i capitolati con almeno ancora un posto libero cioè C1, C3, C4. 
        Inoltre dal documento si spiega come la scelta del gruppo, dopo l'analisi di tutti i capitolati a disposizione, si ricaduta sul capitolato 1: 
        Bot4Me chatbot segretariale, di Imola Informatica.
		\subsection {Riferimenti} 
			\subsubsection {Riferimenti normativi} 
			\begin{itemize}
				\item \textit{Norme di Progetto 0.0.1}
			\end{itemize}
			\subsubsection {Riferimenti informativi} 
			\begin{itemize}
				\item \textbf{Capitolato d'appalto C1 - Bot 4 Me}\newline
				\url{https://www.math.unipd.it/~tullio/IS-1/2021/Progetto/C1.pdf}
				\item \textbf{Verbale Esterno 2022-03-18 con azienda Imola Informatica}
				\item \textbf{Capitolato d'appalto C3 - CC4D}\newline
				\url{https://www.math.unipd.it/~tullio/IS-1/2021/Progetto/C3.pdf}
				\item \textbf{Capitolato d'appalto C4 - Guida Michelin Social}\newline
				\url{https://www.math.unipd.it/~tullio/IS-1/2021/Progetto/C4.pdf}
			\end{itemize}
	\newpage
	%\section{Capitolato C3}
	\begin{itemize}
		\item \textbf{Nome}: \textit{CC4D}: web app di rilevazione conformità;
		\item \textbf{Proponente}: SanMarco Informatica;
		\item \textbf{Committenti}: \textit{Prof. Tullio Vardanega} e \textit{Prof. Riccardo Cardin}.
	\end{itemize}
	\subsection{Descrizione del capitolato}
	Il progetto mira a costruire un'applicazione web che configuri e in seguito visualizzi le \glossario{carte di controllo} delle macchine di un sistema produttivo. Le misurazioni delle carte di controllo sono rilevate e processate da un motore interno che valuta ogni misurazione a seconda dei parametri della configurazione iniziale.
	L'applicativo si suddivide in 4 parti, di seguito spiegate:
	\begin{itemize}
		\item \textbf{web application config}: creare una web application che permetta ad un utente di configurare le macchine del processo con le relative caratteristiche, per ogni caratteristica dev'essere possibile inserire i parametri (media, limiti di controllo) oppure sarà il "motore" che le calcolerà automaticamente relativamente alle ultime N misurazioni
		\item \textbf{API dati}: creare un API che inserisca in un database le misurazioni relative alle caratteristiche delle macchine configurate dall'utente nella web application
		\item \textbf{calculation engine}: creare un motore di calcolo che per ogni misurazione ricevuta la metta in relazione con le precedenti, delle corrette caratteristiche, e la "valuti" per capire se c'è un anomalia o se il processo rimane "in controllo".
		\item \textbf{web application}: creare una web application che permetta all'utente di consultare le caratteristiche configurate per la macchina desiderata e controllare così lo stato del processo per la macchina.
	\end{itemize}
	\subsection{Finalità del progetto}
	L'obiettivo del progetto è quello di mantenere il controllo su un processo produttivo, monitorando costantemente lo stato del processo stesso, mantenendolo in uno stato di "controllo". Le finalità sono, identificare le cause di variabilità nella produzione, diminuire gli "scarti" e continuare costantemente il miglioramento del processo stesso.
	\subsection{Tecnologie interessate}
	Il proponente lascia libera scelta sulle tecnologie da utilizzare, consigliandone però alcune, come spiegato di seguito:
	\begin{itemize}
		\item \textbf{DB Sql o NoSql}: si possono usare entrambi i tipi di database, nel quale si va a memorizzare le configurazioni iniziali dell'utente;
		\item \textbf{time-series DB}: sono database particolari che permettono di gestire serie di dati con un ordine temporale. Da usare per la storcizzazione di tutte le misure rilevate della API;
		\item \textbf{API rest o graphQL}: usati per rilevare le varie misurazioni dalle macchine e trasmetterle al motore dell'applicazione;
		\item \textbf{Java o nodeJS}: per il codice che rigurada il motore dell'applicazione;
		\item \textbf{Angular o React o Vue}: per il codice della web application che espone i dati all'utente;
		\item \textbf{d3js}: è una libreria javascript per manipolare dati organizzati. Da usare per i grafici contenuti nella web application;
	\end{itemize}
	\subsection{Aspetti positivi}
	\begin{itemize}
		\item chiarezza e buona rappresentazione dell'applicativo già nel capitolato;
		\item molto apprezzato perchè può essere di grande aiuto nella realtà aziendale;
	\end{itemize}
	\subsection{Criticità}
	\begin{itemize}
		\item molto complesso con 4 parti che devo interfacciarsi tra di loro;
		\item varie tecnologie diverse, per lo più sconosciute, da integrare nella stessa piattaforma;
	\end{itemize}
	\subsection{Valutazione finale}
	Il capitolato è stato analizzato con attenzione dal gruppo, ed apprezzato per la sua teorica spendibilità nella realtà aziendale, tuttavia però la grande varietà delle tecnologie sconosciute da integrare tutte nella stessa piattaforma, ha spinto il gruppo a non scegliere questo capitolato.
	\newpage
	\section{Capitolato C3}
	\begin{itemize}
		\item \textbf{Nome}: \textit{CC4D}: web app di rilevazione conformità;
		\item \textbf{Proponente}: SanMarco Informatica;
		\item \textbf{Committenti}: \textit{Prof. Tullio Vardanega} e \textit{Prof. Riccardo Cardin}.
	\end{itemize}
	\subsection{Descrizione del capitolato}
	Il progetto mira a costruire un'applicazione web che configuri e in seguito visualizzi le \glossario{carte di controllo} delle macchine di un sistema produttivo. Le misurazioni delle carte di controllo sono rilevate e processate da un motore interno che valuta ogni misurazione a seconda dei parametri della configurazione iniziale.
	L'applicativo si suddivide in 4 parti, di seguito spiegate:
	\begin{itemize}
		\item \textbf{web application config}: creare una web application che permetta ad un utente di configurare le macchine del processo con le relative caratteristiche, per ogni caratteristica dev'essere possibile inserire i parametri (media, limiti di controllo) oppure sarà il "motore" che le calcolerà automaticamente relativamente alle ultime N misurazioni
		\item \textbf{API dati}: creare un API che inserisca in un database le misurazioni relative alle caratteristiche delle macchine configurate dall'utente nella web application
		\item \textbf{calculation engine}: creare un motore di calcolo che per ogni misurazione ricevuta la metta in relazione con le precedenti, delle corrette caratteristiche, e la "valuti" per capire se c'è un anomalia o se il processo rimane "in controllo".
		\item \textbf{web application}: creare una web application che permetta all'utente di consultare le caratteristiche configurate per la macchina desiderata e controllare così lo stato del processo per la macchina.
	\end{itemize}
	\subsection{Finalità del progetto}
	L'obiettivo del progetto è quello di mantenere il controllo su un processo produttivo, monitorando costantemente lo stato del processo stesso, mantenendolo in uno stato di "controllo". Le finalità sono, identificare le cause di variabilità nella produzione, diminuire gli "scarti" e continuare costantemente il miglioramento del processo stesso.
	\subsection{Tecnologie interessate}
	Il proponente lascia libera scelta sulle tecnologie da utilizzare, consigliandone però alcune, come spiegato di seguito:
	\begin{itemize}
		\item \textbf{DB Sql o NoSql}: si possono usare entrambi i tipi di database, nel quale si va a memorizzare le configurazioni iniziali dell'utente;
		\item \textbf{time-series DB}: sono database particolari che permettono di gestire serie di dati con un ordine temporale. Da usare per la storcizzazione di tutte le misure rilevate della API;
		\item \textbf{API rest o graphQL}: usati per rilevare le varie misurazioni dalle macchine e trasmetterle al motore dell'applicazione;
		\item \textbf{Java o nodeJS}: per il codice che rigurada il motore dell'applicazione;
		\item \textbf{Angular o React o Vue}: per il codice della web application che espone i dati all'utente;
		\item \textbf{d3js}: è una libreria javascript per manipolare dati organizzati. Da usare per i grafici contenuti nella web application;
	\end{itemize}
	\subsection{Aspetti positivi}
	\begin{itemize}
		\item chiarezza e buona rappresentazione dell'applicativo già nel capitolato;
		\item molto apprezzato perchè può essere di grande aiuto nella realtà aziendale;
	\end{itemize}
	\subsection{Criticità}
	\begin{itemize}
		\item molto complesso con 4 parti che devo interfacciarsi tra di loro;
		\item varie tecnologie diverse, per lo più sconosciute, da integrare nella stessa piattaforma;
	\end{itemize}
	\subsection{Valutazione finale}
	Il capitolato è stato analizzato con attenzione dal gruppo, ed apprezzato per la sua teorica spendibilità nella realtà aziendale, tuttavia però la grande varietà delle tecnologie sconosciute da integrare tutte nella stessa piattaforma, ha spinto il gruppo a non scegliere questo capitolato.	
	\newpage
	\section{Capitolato C4}
		\begin{itemize}
			\item \textbf{Nome:} \textit{Guida Michelin social}: equivalente di Guida Michelin per post e storie;
			\item \textbf{Proponente}: \textit{Zero12};
			\item \textbf{Committenti}: \textit{Prof. Tullio Vardanega} e \textit{Prof. Riccardo Cardin}.
		\end{itemize}
		\subsection{Descrizione del capitolato}
		Il capitolato proposto dall'azienda \textbf{Zero12} offre la possibilità di realizzare una Guida Michelin rapportata però al mondo dei social con particolare interesse verso le piattaforme di \textbf{Instagram} e \textbf{TikTok}. Incronciando i dati dei post con le recensioni online, sarà possibile creare una mappa di posti suggeriti con la possibilità di specificare persone da seguire per la creazione della guida. La piattoforma fornirà inoltre la possibilità di monitorare le recensioni a partire da un luogo indicato. 
		\subsection{Finalità del progetto}
		Lo scopo di questo progetto è quello di sviluppare un'applicazione mobile \textit{iOS} e \textit{Android} in grado di mostrare la \textit{Guida Social Michelin}. Per realizzare questo obiettivo il team dovrebbe svolgere un'analisi sulle API di TikTok e Instagram al fine di indentificare il miglior approccio per la raccolta e l'analisi delle informazioni. Basandosi successivamente su un'architettura a microservizi per la creazione effettiva dell'applicativo. 
		\subsection{Tecnologie interessate}
		Per la realizzazione di questo capitolato, l'azienda proponente consiglia di utilizzare gli \textbf{Amazon Web Services}. In particolare facendo riferimento ai servizi di:
			\begin{itemize}
				\item \textbf{AWS Fargate}: servizio che permette la creazione di applicativi senza dover gestire i server
				\item \textbf{AWS Appsync}: servizio completamente gestito che facilita lo sviluppo di API GraphQL
				\item \textbf{Amazon Neptune}: servizio di database a grafo
			\end{itemize}
		Per quanto riguarda invece i linguaggi di programmazione, sarebbero necessari: \textbf{NodeJS}, \textbf{Swift} e \textbf{Kotlin}. Il primo utilizzato per sviluppare API Restful JSON, mentre i restanti verrebbero utilizzati rispettivamente per programmare applicazioni in ambito iOS e Android. 
		\subsection{Aspetti positivi}
		Il progetto proposto ha diversi aspetti positivi, in particolare la possibilità offerta dall'azienda \textbf{Zero12} di usufruire di attività di formazione sia sulle tecnologie AWS che su quelle Mobile che permetterebbe di coprire alcune lacune presenti su tali argomenti. Inoltre il gruppo ritiene che la possibilità di lavorare sull'elaborazione dei dati di due piattaforme social così importanti costituisca un ottimo bagaglio culturale da poter apprendere nel corso di questa offerta formativa. 
		\subsection{Criticità}
		Nello specifico il gruppo non ha rilevato particolari criticità relative allo sviluppo del progetto proposto. Le uniche incertezze emerse riguardano il livello di difficoltà nel riuscire a manipolare i dati offerti dalle piattoforme \textbf{Instagram} e \textbf{TikTok} al fine di ottenere il risultato voluto. Tuttavia non si ritiene essere un elemento limitante per lo sviluppo del progetto grazie anche al supporto eventuale offerto dall'azienda \textbf{Zero12}.
		\subsection{Valutazione finale}
		Nonostante il capitolato proposto risulti essere molto interessante offrendo la possibilità di lavorare con delle tecnologie che fornirebbero un bagaglio di cononscenze molto utile, a seguito di una votazione interna il team ha deciso di concentrare la propria candidatura su un altro progetto. 
	\newpage
	\section{Conclusioni}
	    \subsection {} Inserire qui le conclusioni
	\newpage
	
	\section*{Tracciamento delle decisioni}
	\renewcommand{\arraystretch}{1.8} % Aumento ampiezza righe
	\begin{tabular}{ |c|l| }
		\hline
		\textbf{Codice} & \textbf{Descrizione} \\
		\hline
		VI\_aaaa-mm-gg.1 & Riempire qui se verbale interno\\ \hline
		VE\_aaaa-mm-gg.1 & Riempire qui se verbale esterno\\ \hline
		% Inserire altre righe per nuovi cambiamenti
	\end{tabular}
\end{document}